\divider{Higher Sobolev Continuation}

Enstrophy is the first derivative level at which the energy estimate loses pointwise-in-time control.
At higher Sobolev orders, the same first spatial derivative reappears in a stronger role: its \(L^\infty\) size controls the growth of every higher derivative.
For a fixed integer \(m\geq3\), apply every spatial derivative through order \(m\), pair the resulting equations with the same derivatives of \(u\), and sum.
For each differentiated row, the pressure pairing vanishes because \(\nabla\cdot\partial^\alpha u=0\).
The undifferentiated transport also cancels, while the remaining transport commutators place one spatial derivative on the carrying field.
The resulting estimate has the form
\[
  \frac12\frac{d}{dt}\lVert u(t)\rVert_{H^m(\mathbb R^3)}^2
  +\nu\lVert\nabla u(t)\rVert_{H^m(\mathbb R^3)}^2
  \leq
  C_m\lVert\nabla u(t)\rVert_{L^\infty(\mathbb R^3)}
  \lVert u(t)\rVert_{H^m(\mathbb R^3)}^2.
\]
The \(L^\infty\) norm of \(\nabla u\) is the coefficient through which transport can amplify every differentiated row.
If its accumulated size is finite,
\[
  \int_0^T
  \lVert\nabla u(t)\rVert_{L^\infty(\mathbb R^3)}\,dt
  <\infty,
\]
then Gr\"onwall's inequality keeps \(\lVert u(t)\rVert_{H^m}\) finite through \(T\).
The initial datum is smooth, so every finite \(H^m\) norm is present at time zero.
Each finite order carries its own constant \(C_m\), so the argument asks for separate finite control at each \(m\).
One accumulated deformation bound propagates each finite spatial order separately.

Now choose one classical continuation level \(s>5/2\).
Sobolev embedding makes a uniform \(H^s\) bound a uniform bound on \(\nabla u\).
Local well-posedness in \(H^s(\mathbb R^3)\) also gives a common forward lifespan for every initial state in a bounded \(H^s\) ball \parencite{FujitaKato1964,Kato1984}.
Suppose the smooth solution exists on \([0,T)\) and
\[
  \sup_{0\leq t<T}
  \lVert u(t)\rVert_{H^s(\mathbb R^3)}
  \leq M.
\]
The local lifespan has a positive lower bound \(\tau=\tau(s,M,\nu)\).
Choose \(t_0<T\) with \(T-t_0<\tau\), and restart local theory from the actual state \(u(t_0)\).
The restarted solution exists past \(T\), and strong uniqueness makes it agree with the original solution wherever both are defined.
The restart begins from an actual preterminal state \(u(t_0)\), and the extended solution itself supplies the value at \(T\).

Continuation now fixes the target for a global argument.
Finiteness on each strict subinterval still allows growth as the candidate endpoint approaches; continuation needs control that remains uniform across those expanding intervals.
For the smoothness required here, every finite spatial order of \(u\), together with the pressure and temporal responses generated by the equation, must belong to that same history at the terminal face.
Once that terminal state is available, local existence and uniqueness carry the history forward.

A finite maximal time \(T_*<\infty\) would force every continuation-grade Sobolev norm to become unbounded and
\[
  \int_0^{T_*}
  \lVert\nabla u(t)\rVert_{L^\infty(\mathbb R^3)}\,dt
  =\infty.
\]
Leray's energy inequality stops at \(L^\infty_tL^2_x\cap L^2_tH^1_x\), below both restart controls.
Caffarelli, Kohn, and Nirenberg next asked what the local energy inequality can still force near any spacetime point where continuation might fail.
